\documentclass[11pt]{article}
\usepackage[a4paper,margin=25mm]{geometry}
\usepackage{amsmath,amssymb,amsthm,booktabs}
\usepackage[hidelinks]{hyperref}
\emergencystretch=2em
\newtheorem{theorem}{Theorem}
\newtheorem{lemma}{Lemma}
\newcommand{\Z}{\mathbb Z}
\title{Equal-size generating sets give\\different Cayley-core edge counts\\
\large Disproof of Conjecture 00000003467}
\author{gaochengzhecpu}\date{}
\begin{document}\maketitle
\begin{abstract}
On the same cyclic group $\Z/12\Z$, the inverse-closed generating sets
$S_1=\{\pm1,\pm3\}$ and $S_2=\{\pm1,\pm2\}$ both have size $4$.
Their connected Cayley graphs have cores $K_2$ and $K_3$, respectively,
with $1$ and $3$ edges. Consequently the core edge count is not a
function of generating-set size even when the group is fixed.
This disproves the asserted linear dependence in Conjecture 00000003467,
independently of its specified slope.
\end{abstract}

\section*{Definitions and scope}
We use finite simple undirected graphs. For an additive group $G$ and
a subset $S=-S$ not containing $0$, the Cayley graph has vertex set
$G$ and adjacency $x\sim y$ if $y-x\in S$.
A graph homomorphism preserves adjacency. A core of a finite graph is
a vertex-minimal subgraph to which the graph has a homomorphism.
Equivalently, it is a retract whose endomorphisms are injective.
The standard Cayley and core conventions are described in Roberson,
Introduction and Section 1. We give the necessary correspondence below.

The source asserts a linear dependence on generating-set size for
abelian groups. We keep the entire group fixed, so its rank and every
other group invariant remain fixed. Both connection sets below genuinely
generate the group and the graphs are connected. The source does not
require generating sets to be irredundant.

\section*{Two Cayley graphs on the same group}
Let $G=\Z/12\Z$ and put
\[
 S_1=\{1,11,3,9\}=\{\pm1,\pm3\},\qquad
 S_2=\{1,11,2,10\}=\{\pm1,\pm2\}.
\]
Both sets have four elements, omit $0$, and are closed under negation.
Each contains $1$, which generates $G$. In either graph, repeatedly
following the edge $x\sim x+1$ reaches every vertex, proving connectedness.
Write $A=\operatorname{Cay}(G,S_1)$ and
$B=\operatorname{Cay}(G,S_2)$.

For $A$, take the induced edge on vertices $\{0,1\}$ and define
\[
 r_2(x)=x\bmod2\in\{0,1\}.
\]
This is well-defined on $\Z/12\Z$ because $2$ divides $12$.
Every difference in $S_1$ is odd, so adjacent vertices have distinct
images. Thus $r_2$ is a graph homomorphism to this copy of $K_2$.
It fixes $0$ and $1$, making it a retraction.

For $B$, the vertices $\{0,1,2\}$ form an induced triangle: all their
nonzero differences lie in $S_2$. Define
\[
 r_3(x)=x\bmod3\in\{0,1,2\}.
\]
This is well-defined since $3$ divides $12$. None of $1,11,2,10$ is
zero modulo $3$, so $r_3$ preserves adjacency to $K_3$.
It fixes the three triangle vertices and is therefore a retraction.

\section*{Why these retracts are the actual cores}
\begin{lemma}
If a finite graph $X$ retracts onto a complete graph $K_t$, then $K_t$
is its core, unique up to graph isomorphism.
\end{lemma}
\begin{proof}
Every homomorphism out of $K_t$ is injective: identifying two vertices
would turn their edge into a loop. If $X$ maps to a subgraph $Y$, the
restriction to the displayed $K_t$ therefore forces $|V(Y)|\ge t$.
Since the retraction onto $K_t$ attains this bound, it is a core.
For a vertex-minimal choice $Y$, the same restriction is an injection
from $K_t$ onto its $t$ vertices. It preserves every edge of $K_t$,
so $Y$ is complete and isomorphic to $K_t$.
\end{proof}

Applying the lemma to the explicit retractions yields
\[
 \operatorname{core}(A)\cong K_2,\qquad
 \operatorname{core}(B)\cong K_3.
\]
In particular their edge counts are exactly $1$ and $3$; they are not
merely upper bounds inferred from colorings.

\begin{theorem}
The abelian core-edge-count assertion of Conjecture 00000003467 is false.
\end{theorem}
\begin{proof}
If a function $F_G$ of generating-set size determined the core edge
count for this fixed group, the two examples would give
\[
          F_G(4)=1,\qquad F_G(4)=3,
\]
a contradiction. Thus no such function exists, in particular no
linear or affine formula with slope half the group rank. Failure of
the abelian assertion already disproves the conjunction; no statement
about the separate nonabelian bound is needed.
\end{proof}

\section*{Formal correspondence}
The Lean project constructs the actual ring \texttt{ZMod 12} and two
\texttt{SimpleGraph} objects directly from the difference-membership
Cayley definition. It checks the generating-set sizes, proves their
additive subgroup closures are the whole group, and proves graph
connectedness by walks of unit steps. The quotient-color maps and
inclusions are actual graph homomorphisms, and their left-inverse
identities are checked.

The local predicate \texttt{IsCore} is the standard injective-endomorphism
condition; \texttt{IsCoreOf} adds an explicit retraction.
The theorem \texttt{retraction\_reflects\_adjacency} checks that the
inclusion identifies an induced subgraph.
\texttt{core\_minimal} proves the minimum-vertex property against any
finite subgraph candidate, so the core predicate is connected to the
usual definition rather than merely naming a convenient retract.

The formal uniqueness proof also covers every finite core satisfying
this standard predicate. If $H$ is another core and there are
homomorphisms $f:H\to K_t$ and $g:K_t\to H$, the endomorphism $gf$
is injective, hence surjective by finiteness. Consequently $g$ is
surjective; it is also injective because its source is complete.
It therefore identifies $H$ with $K_t$ as a graph. This is proved in
\texttt{complete\_core\_unique}, and
\texttt{complete\_core\_edge\_count} transfers the actual finite edge
count to any such core. The final counterexample and
\texttt{no\_function\_of\_generator\_count} use these verified cores
and Mathlib's unordered edge sets.

\section*{Reproduction}
Lean 4.19.0 and Mathlib revision\\
\texttt{c44e0c8ee63ca166450922a373c7409c5d26b00b}
are pinned. In \texttt{lean/}, run
\begin{verbatim}
lake exe cache get
lake build
lake env lean -DwarningAsError=true Main.lean
\end{verbatim}
The cache command fetches pinned official dependencies if absent.
The submitted project itself is built in a fresh directory, with
warnings treated as errors. Finite calculations are kernel checked;
no external numerical oracle, extra axiom or proof placeholder is used.
No auxiliary numerical program is required. Run
\texttt{tectonic main.tex} to regenerate the PDF.
The author performed a separate adversarial self-review; no independent
reviewer or subagent was used. Exact source, provenance, build logs,
axiom dependencies, hashes and review records accompany this submission.

\section*{References}
\begin{enumerate}
\item The Last Math Competition, Conjecture 00000003467.
The exact bilingual source is reproduced in \texttt{SOURCE.md}.
\item D. E. Roberson, \emph{Cores of Vertex Transitive Graphs},
Electronic Journal of Combinatorics 20(2) (2013), P45,
Introduction and Section 1.
\href{https://www.combinatorics.org/ojs/index.php/eljc/article/download/v20i2p45/pdf/}{Publisher's full text}.
\end{enumerate}
\end{document}
